\section{Numerical Experiments}
\label{sec:sims}

All experiments use the simulation model \eqref{eq:M} with $\zeta = 1$,
$\sigma_0 = 1$, $H = 0.6$ (unless stated), and the two-step FWLS procedure
(\Cref{alg:fwls}).  We used 8000 Monte Carlo replications for size and
5000 for power; detailed numerical tables are reported in
Appendix~\ref{app:numerical_tables}.

\subsection{Null distribution calibration}

Figure~\ref{fig:null-calibration} shows that the residual adequacy statistic
is already well calibrated at moderate lag counts: its empirical mean tracks
the oracle null mean closely, and the rejection frequency remains near the
nominal level.  Table~\ref{tab:null} gives the exact grid values.

\begin{figure}[!htbp]
\centering
\safeincludegraphics[width=\linewidth]{artifacts/figures/scale_consistency/fig_null_calibration}
\caption{Null calibration of the residual adequacy statistic.}
\label{fig:null-calibration}
\end{figure}
\FloatBarrier

\subsection{Power at the minimax boundary}

Figure~\ref{fig:power-boundary} shows the transition at the minimax boundary:
the adequacy procedure is informative but not overconfident at $c=1$, and it
becomes decisive above the boundary.  Table~\ref{tab:power-boundary} records
the exact values.

\begin{figure}[!htbp]
\centering
\safeincludegraphics[width=\linewidth]{artifacts/figures/scale_consistency/fig_power_boundary}
\caption{Power transition at the minimax boundary.}
\label{fig:power-boundary}
\end{figure}
\FloatBarrier

\subsection{FWLS versus oracle}

Figure~\ref{fig:fwls-oracle-gap} shows that the feasible estimator tracks the
oracle benchmark closely: the gap shrinks with $n$ and remains small across
the tested lag lengths.  This is the numerical counterpart of the CLT.

\begin{figure}[!htbp]
\centering
\safeincludegraphics[width=\linewidth]{artifacts/figures/scale_consistency/fig_fwls_oracle_gap}
\caption{Agreement between feasible and oracle estimation.}
\label{fig:fwls-oracle-gap}
\end{figure}
\FloatBarrier

\subsection{Rate constant}

Figure~\ref{fig:rate-constant} shows that the FWLS error obeys the predicted
scaling over two orders of magnitude in $n$, while the rescaled constant
remains essentially flat.  Table~\ref{tab:rate} gives the exact grid values.

\begin{figure}[!htbp]
\centering
\safeincludegraphics[width=\linewidth]{artifacts/figures/scale_consistency/fig_rate_constant}
\caption{Rate-constant stabilization for $\widehat H$.}
\label{fig:rate-constant}
\end{figure}
\FloatBarrier

\subsection{Robustness}

Figure~\ref{fig:robustness} varies the scale law on the left and the noise law
on the right.  The procedure is stable under moderate perturbations and
degrades visibly only when the structural assumptions are pushed beyond the
scope of the theorem.

\begin{figure}[!htbp]
\centering
\safeincludegraphics[width=\linewidth]{artifacts/figures/scale_consistency/fig_robustness}
\caption{Sensitivity to scale misspecification and alternative noise laws.}
\label{fig:robustness}
\end{figure}
\FloatBarrier

Tables~\ref{tab:null}, \ref{tab:power-boundary}, and \ref{tab:rate}, together
with the appendix robustness tables,
provide the exact numerical grid values.

\FloatBarrier
